% ~5,5 páginas — é a secção principal do artigo.
\section{Estado da arte}
\label{sec:estado-arte}

% ---------------------------------------------------------------
\subsection{Dados e métricas de previsão}
\label{sec:dados-metricas}

A investigação em previsão de movimento apoia-se em conjuntos de dados
públicos. O WOMD~\cite{ettinger2021large} tem mais de 100\,000 cenas de
\un{20}{s}, num total de 574 horas a \un{10}{Hz}, com as trajetórias de
veículos, peões e ciclistas e o mapa da zona. A tarefa oficial dá
\un{1}{s} de passado e pede os \un{8}{s} seguintes. O
Argoverse~\cite{chang2019argoverse} tem 324\,557 sequências de \un{5}{s},
com \un{2}{s} observados e \un{3}{s} a prever, e distingue-se por
publicar, ao lado dos modelos neuronais, referências que não o são:
velocidade constante e regressão por vizinhos mais próximos.

As métricas habituais medem distâncias entre a trajetória prevista e a
real~\cite{chang2019argoverse,ettinger2021large}. O erro médio de
deslocamento (ADE) é a média dessa distância ao longo do horizonte; o
erro final (FDE) é a distância no último instante. Quando o modelo propõe
$K$ trajetórias, conta a melhor delas (minADE, minFDE). Neste trabalho
prevê-se uma única trajetória por veículo, pelo que a métrica é o FDE.

% ---------------------------------------------------------------
\subsection{Previsão de movimento}
\label{sec:previsao}

A classificação mais citada separa os modelos em três níveis: os físicos,
que só usam a dinâmica do veículo; os baseados em manobras, que primeiro
reconhecem a intenção do condutor; e os que consideram a interação entre
veículos~\cite{lefevre2014survey}. Uma revisão mais recente organiza os
métodos pela técnica: modelos físicos, aprendizagem automática clássica,
aprendizagem profunda e aprendizagem por
reforço~\cite{huang2022survey}. É a segunda família que este trabalho
explora.

\paragraph{Modelos físicos.}
Assumem que o veículo mantém a velocidade, a aceleração ou a taxa de
viragem que tem no momento. São os modelos usados no seguimento de
veículos com filtros de Kalman, contexto em que foram
comparados~\cite{schubert2008comparison}. Não precisam de treino,
mas só são fiáveis a muito curto prazo, abaixo de um
segundo~\cite{lefevre2014survey}: ignoram a estrada, os semáforos e os
outros veículos.

Ainda assim, são referências difíceis de bater, e por isso qualquer
modelo aprendido deve ser comparado com elas. Na previsão de peões, a
velocidade constante superou modelos neuronais do estado da arte nos
conjuntos de dados habituais~\cite{scholler2020constant}. Para veículos,
a \cref{tab:baselines} reúne valores publicados. A distância para os
modelos aprendidos cresce com o horizonte, e no WOMD a velocidade
constante erra mais de \un{5}{m} logo aos \un{3}{s}.

\begin{table}[t]
\centering
\caption{Erros publicados da velocidade constante e de modelos aprendidos
  na previsão de veículos, em metros, com uma trajetória proposta por
  veículo, exceto onde se indica $K$.}
\label{tab:baselines}
\begin{tabular*}{\textwidth}{@{\extracolsep{\fill}}llllll@{}}
\toprule
Dados & Horizonte & Métrica & Vel.\ const. & Modelo aprendido & Fonte \\
\midrule
Argoverse & \un{3}{s} & minFDE & 7,89 & 4,97 (LSTM)
  & \cite{chang2019argoverse} \\
NGSIM & \un{3}{s} / \un{5}{s} & RMSE & 3,13 / 6,68
  & 2,09 / 4,37 (CS-LSTM) & \cite{deo2018convolutional} \\
WOMD & \un{3,1}{s} / \un{5,1}{s} & minFDE & 5,7 / 14,1 & --
  & \cite{munozsanchez2022scenario} \\
WOMD & \un{8}{s} & minADE & 11,0 & 1,34 (LSTM, $K{=}6$)
  & \cite{ettinger2021large} \\
\bottomrule
\end{tabular*}
\end{table}

\paragraph{Aprendizagem automática clássica.}
Antes das redes profundas, a previsão fazia-se com classificadores e
regressores sobre variáveis calculadas à mão. Em autoestrada, uma floresta aleatória
estima a probabilidade de cada manobra e uma regressão por misturas
gaussianas converte-a numa distribuição de
trajetórias~\cite{schlechtriemen2015when}. O estudo mais sistemático
nesta linha, com mais de 30\,000~km de dados de autoestrada, concluiu que
a floresta aleatória classifica as manobras melhor do que um perceptrão
multicamada e do que um classificador bayesiano ingénuo, e obteve um erro
lateral mediano inferior a \un{0,21}{m} a
\un{5}{s}~\cite{wirthmuller2021teaching}. O Argoverse dá o contraponto:
a regressão por vizinhos mais próximos, aplicada diretamente às
trajetórias passadas, fica ao nível da velocidade constante
(\un{7,88}{m} contra \un{7,89}{m} a
\un{3}{s})~\cite{chang2019argoverse}. O modelo, por si só, não chega; contam
as variáveis que recebe.

As florestas aleatórias~\cite{breiman2001random} e o \emph{gradient
boosting}~\cite{friedman2001greedy} são hoje os modelos de referência
para dados em tabela. Estes trabalhos concentram-se, porém, em classificar
manobras em contextos específicos, sobretudo a mudança de faixa em
autoestrada. Na pesquisa feita para este artigo não se encontrou nenhum
que aplique modelos clássicos à regressão de trajetórias no WOMD, em meio
urbano.

\paragraph{Aprendizagem profunda.}
Os melhores resultados atuais vêm de redes neuronais que recebem as
trajetórias e o mapa quase em bruto. A VectorNet representa faixas e
agentes como vetores ligados num grafo~\cite{gao2020vectornet}; o Motion
Transformer, que liderou a tabela do WOMD, atinge um minFDE de
\un{1,22}{m} com seis trajetórias propostas, em média sobre os horizontes
de 3, 5 e \un{8}{s}~\cite{shi2022motion}. Estes números não se comparam
diretamente com os de uma previsão única. O preço é o volume de dados, o
treino em GPU e a dificuldade em explicar cada previsão.

A \cref{tab:familias} resume as três famílias.

\begin{table}[t]
\centering
\caption{Famílias de modelos de previsão de movimento.}
\label{tab:familias}
\begin{tabular*}{\textwidth}{@{\extracolsep{\fill}}llll@{}}
\toprule
Família & Entrada & Treino & Explicação da previsão \\
\midrule
Física~\cite{schubert2008comparison} & estado atual & nenhum & total (uma fórmula) \\
Clássica~\cite{wirthmuller2021teaching} & variáveis em tabela & CPU, poucos dados
  & peso de cada variável \\
Profunda~\cite{gao2020vectornet,shi2022motion} & trajetórias e mapa & GPU, muitos dados
  & limitada \\
\bottomrule
\end{tabular*}
\end{table}

\paragraph{Duas ideias que não dependem do modelo.}
A primeira é o referencial. Os modelos centrados no agente, que descrevem
a cena a partir da posição e da orientação do veículo a prever, são
invariantes a translações e rotações e dominam as tabelas
públicas~\cite{su2022narrowing,shi2022motion}. A segunda é aprender uma
correção em vez da trajetória inteira: parte-se da previsão de um modelo
baseado em conhecimento e o modelo aprendido estima apenas o resíduo, o
que melhora a exatidão e a generalização a cenas
novas~\cite{bahari2021injecting}. Ambas se aplicam a modelos clássicos e
são a base da abordagem da \cref{sec:abordagem}.

% ---------------------------------------------------------------
\subsection{Planeamento de movimento}
\label{sec:planeamento}

O planeamento divide-se habitualmente em níveis: a rota na rede viária, a
decisão de comportamento (seguir, ceder, mudar de faixa), o planeamento de
movimento, que produz uma trajetória, e o
controlo~\cite{paden2016survey}. No planeamento de movimento há três
grandes famílias: pesquisa em grafos, amostragem e otimização
numérica~\cite{paden2016survey}.

\paragraph{Pesquisa em grafos.}
O espaço de estados é discretizado e procura-se o caminho de menor custo.
O A* expande primeiro os estados com menor soma do custo já gasto com uma
estimativa do custo em falta; se essa heurística nunca sobrestimar, a
solução é ótima~\cite{hart1968formal}. Para lidar com obstáculos móveis, o
tempo e a velocidade passam a fazer parte do estado. O trabalho mais
próximo do aqui
proposto pesquisa sobre o tempo, a distância percorrida, a faixa e a
velocidade discreta, trata os outros veículos como obstáculos no
espaço-tempo, previstos com velocidade constante e uma margem de
segurança, e usa como heurística o custo de um problema
relaxado~\cite{ajanovic2018search}.

\paragraph{Decomposição caminho-velocidade.}
Uma simplificação clássica resolve o problema em dois passos: primeiro
um caminho que evita os obstáculos fixos, depois um perfil de velocidade
ao longo dele que evita os obstáculos móveis no plano distância-tempo
(S--T)~\cite{kant1986toward}. Em estrada, o caminho natural é o eixo da
faixa, e as trajetórias descrevem-se em coordenadas de Frenet (distância
ao longo da faixa e desvio lateral), onde se geram polinómios que
minimizam a variação da aceleração~\cite{werling2010optimal}. O planeador
EM do Apollo, de uso industrial, segue esta decomposição: otimiza o
caminho e a velocidade em separado, com programação dinâmica para uma
solução grosseira e programação quadrática para a
refinar~\cite{fan2018baidu}.

\paragraph{Modelos de seguimento.}
O Intelligent Driver Model (IDM) calcula a aceleração a partir da
distância e da velocidade relativa ao veículo da
frente~\cite{treiber2000congested}; é o modelo que controla os outros
veículos nos simuladores
atuais~\cite{karnchanachari2024towards,gulino2023waymax}.

\paragraph{Porquê pesquisa neste trabalho.}
Com ações discretas (acelerar, manter, travar) e obstáculos que ocupam
zonas do plano distância-tempo, o problema é naturalmente um grafo. A
pesquisa é determinística: a mesma entrada dá sempre o mesmo plano, ao
contrário dos métodos por amostragem, que são aleatórios por
construção~\cite{karaman2011sampling}. E não precisa de uma solução
inicial, que os métodos por otimização exigem e que o próprio planeador
do Apollo vai buscar a uma pesquisa~\cite{fan2018baidu}.

% ---------------------------------------------------------------
\subsection{Da previsão ao planeamento}
\label{sec:integracao}

\paragraph{Formas de integração.}
Uma revisão recente distingue três: sequencial, em que a previsão é feita
primeiro e entra no planeador como dado de entrada ou termo de custo; não
dirigida, em que as duas tarefas são resolvidas em conjunto; e
bidirecional, em que o plano influencia a
previsão~\cite{hagedorn2025integration}. A abordagem deste trabalho é
sequencial. A sua limitação conhecida é tratar os outros como obstáculos
que não reagem, o que torna o veículo defensivo; modelar o efeito das
ações próprias nos outros condutores dá comportamentos mais
eficientes~\cite{sadigh2016planning}.

\paragraph{Uma previsão melhor dá um plano melhor?}
Não necessariamente. A exatidão medida em dados fixos não acompanha de
forma fiável o desempenho na condução, por duas razões: o plano do veículo
altera o comportamento dos outros, e um preditor mais pesado atrasa o
planeador~\cite{tran2023what}. O exemplo
mais forte é o PDM-Closed, um planeador por regras que prevê os outros
veículos com velocidade constante e escolhe entre propostas geradas pelo
IDM: no nuPlan obteve 92 e 93 pontos em simulação, contra 72 e 73 do
melhor planeador aprendido (com e sem veículos reativos), apesar de ter
a pior pontuação quando avaliado apenas contra o
registo~\cite{dauner2023parting}.

\paragraph{Como se avalia um planeador.}
Em ciclo aberto, compara-se o plano com o que o condutor real fez. Em
ciclo fechado, o plano é executado num simulador e o planeador volta a
decidir a cada passo. Os outros veículos podem repetir o registo, sem
reagir, ou ser controlados por um modelo reativo como o
IDM~\cite{karnchanachari2024towards}. O nuPlan mede colisões por culpa
própria, permanência na zona transitável, progresso, tempo até à colisão e
conforto~\cite{karnchanachari2024towards}; o Waymax oferece as mesmas
duas opções sobre as cenas do WOMD~\cite{gulino2023waymax}. A repetição
do registo tem um defeito conhecido: quando o plano se afasta do que
aconteceu, a simulação diverge e surgem colisões com veículos que, na
realidade, teriam reagido~\cite{karnchanachari2024towards,hagedorn2025integration}.

% ---------------------------------------------------------------
\subsection{Síntese}
\label{sec:sintese}

Da revisão saem três conclusões que orientam o resto do trabalho.
Primeiro, a velocidade constante é a referência obrigatória e o seu erro
no WOMD é conhecido (\cref{tab:baselines}). Segundo, os modelos clássicos
têm resultados sólidos na classificação de manobras em autoestrada, mas
não foram encontrados trabalhos que os usem para prever trajetórias no
WOMD. Terceiro, o ganho em exatidão só interessa se chegar ao plano, e
isso mede-se com o planeador, não com o erro de previsão. Tanto quanto
foi possível apurar, não existe uma comparação, com um planeador por
pesquisa sobre estes dados, entre quatro fontes de previsão: nenhuma,
velocidade constante, modelo clássico e futuro real.
